\section{Additional Tables and Figures}
\label{app:additional-descriptives}

This appendix reports the detailed employment, income, wage-growth, and
worker-characteristics tables that support the compact exhibits in the main
text. It also reports median real hourly wages by firm size and worker group for
2008 and 2025, 2025 wage gaps by worker group, plus supporting descriptive
figures and regression robustness exercises.

\subsection{Detailed Descriptive Tables}

Table~\ref{tab:descriptive-employment-size} reports the pooled employment
distribution by firm size. Tables~\ref{tab:descriptive-income-size-comparison}
and \ref{tab:descriptive-worker-characteristics-comparison} compare the pooled
sample with 2025. Tables~\ref{tab:descriptive-income-size} and
\ref{tab:descriptive-worker-characteristics} report the pooled income and worker
composition profiles by firm size in more detail. Tables
\ref{tab:descriptive-income-size-2025} and
\ref{tab:descriptive-worker-characteristics-2025} repeat the same exercise for
2025. Table~\ref{tab:descriptive-wage-growth-2008-2025} reports wage changes
between 2008 and 2025, Table~\ref{tab:descriptive-wage-groups-mean-2008-2025}
reports the corresponding mean wages by worker group, Table~\ref{tab:descriptive-wage-groups-median-2008-2025}
reports medians, and
Table~\ref{tab:descriptive-wage-gaps-2025} reports 2025 wage gaps by sex,
formality status, education group, and sector.
The appendix tables reinforce the same message as the main text. The raw
firm-size premium is visible in means, quartiles, and medians; the 2025-only
tables look similar to the pooled tables; and the composition tables show that
larger firms are simultaneously more formal and more intensive in higher
education. The wage-gap table adds that formality, education, and sector
differences are large even before conditioning on firm size, so observable
sorting is an important part of the empirical environment.

\begin{table}[p]
  \centering
  \caption{Mean real hourly labor income gaps by worker group, 2025}
  \label{tab:descriptive-wage-gaps-2025}
  \scriptsize
  \begin{tabular}{lp{0.38\textwidth}rrr}
    \toprule
    Dimension & Group & Employment share & Mean income & Gap \\
    \midrule
    Sex & Men & 54.3\% & 12,849 & 1.6\% \\
    Sex & Women & 45.7\% & 12,408 & -1.9\% \\
    Formality & Formal & 57.2\% & 16,571 & 31.0\% \\
    Formality & Informal & 42.8\% & 7,401 & -41.5\% \\
    Education & Preschool or less & 0.9\% & 4,497 & -64.4\% \\
    Education & Primary & 10.6\% & 6,282 & -50.3\% \\
    Education & Secondary & 44.3\% & 7,751 & -38.7\% \\
    Education & Higher education & 44.3\% & 19,214 & 51.9\% \\
    Sector & T - Household employers & 2.9\% & 7,183 & -43.2\% \\
    Sector & I - Accommodation and food & 7.8\% & 7,737 & -38.8\% \\
    Sector & H - Transport & 8.8\% & 8,524 & -32.6\% \\
    Sector & F - Construction & 6.9\% & 9,385 & -25.8\% \\
    Sector & G - Commerce & 19.4\% & 9,808 & -22.4\% \\
    Sector & R/S - Arts and other services & 6.1\% & 9,999 & -20.9\% \\
    Sector & C - Manufacturing & 12.8\% & 11,010 & -12.9\% \\
    Sector & D/E - Utilities and water & 1.6\% & 12,795 & 1.2\% \\
    Sector & A - Agriculture & 0.9\% & 13,975 & 10.5\% \\
    Sector & L/M/N - Business services & 12.2\% & 15,104 & 19.4\% \\
    Sector & Q - Health & 5.6\% & 16,668 & 31.8\% \\
    Sector & P - Education & 5.0\% & 20,772 & 64.2\% \\
    Sector & K - Finance and insurance & 2.7\% & 22,241 & 75.9\% \\
    Sector & J - Information and communication & 2.6\% & 23,057 & 82.3\% \\
    Sector & O - Public administration & 4.4\% & 23,893 & 88.9\% \\
    Sector & B - Mining & 0.3\% & 30,685 & 142.6\% \\
    \bottomrule
  \end{tabular}
  \vspace{0.3em}
  \begin{minipage}{0.95\textwidth}
  \footnotesize
  Notes: Statistics use 2025 observations and GEIH expansion weights. Employment shares are computed within each dimension. Mean income is real hourly labor income in constant 2025 pesos. The gap is the percent difference in each group's weighted mean relative to the overall 2025 weighted mean. Formality rows exclude pensioned occupied workers. Sector rows omit extraterritorial organizations.
  \end{minipage}
\end{table}



\begin{table}[htbp]
  \centering
  \caption{Employment distribution by firm size}
  \label{tab:descriptive-employment-size}
  \small
  \begin{tabular}{lrrrrr}
    \toprule
    Firm size & Observations & Worker-years (m) & Pooled share & 2008 share & 2025 share \\
    \midrule
    Solo & 1,435,259 & 50.6 & 31.1\% & 34.5\% & 28.9\% \\
    2-3 & 494,122 & 19.6 & 12.0\% & 16.6\% & 9.8\% \\
    4-5 & 236,560 & 9.9 & 6.1\% & 6.8\% & 6.2\% \\
    6-10 & 213,420 & 10.1 & 6.2\% & 5.9\% & 6.3\% \\
    11-19 & 141,540 & 6.8 & 4.2\% & 3.7\% & 4.6\% \\
    20-30 & 128,837 & 6.5 & 4.0\% & 3.7\% & 4.1\% \\
    31-50 & 105,125 & 5.5 & 3.4\% & 3.0\% & 3.8\% \\
    51-100 & 111,905 & 6.1 & 3.8\% & 3.5\% & 4.4\% \\
    101+ & 985,554 & 47.5 & 29.2\% & 22.2\% & 31.9\% \\
    \bottomrule
  \end{tabular}
  \vspace{0.3em}
  \begin{minipage}{0.95\textwidth}
  \footnotesize
  Notes: Worker-years and employment shares use GEIH expansion weights. The sample matches the baseline regression sample.
  \end{minipage}
\end{table}


\begin{table}[htbp]
  \centering
  \caption{Real hourly labor income by firm size: pooled sample and 2025}
  \label{tab:descriptive-income-size-comparison}
  \small
  \begin{tabular}{lrrrr}
    \toprule
    Firm size & Pooled mean & 2025 mean & Pooled premium & 2025 premium \\
    \midrule
    Solo & 7,248 & 8,351 & 0.0\% & 0.0\% \\
    2-3 & 7,959 & 8,830 & 9.8\% & 5.7\% \\
    4-5 & 8,912 & 9,503 & 23.0\% & 13.8\% \\
    6-10 & 9,608 & 10,609 & 32.6\% & 27.0\% \\
    11-19 & 10,852 & 12,346 & 49.7\% & 47.8\% \\
    20-30 & 10,782 & 11,997 & 48.8\% & 43.6\% \\
    31-50 & 11,576 & 13,428 & 59.7\% & 60.8\% \\
    51-100 & 12,697 & 14,481 & 75.2\% & 73.4\% \\
    101+ & 15,735 & 18,502 & 117.1\% & 121.5\% \\
    \bottomrule
  \end{tabular}
  \vspace{0.3em}
  \begin{minipage}{0.95\textwidth}
  \footnotesize
  Notes: Pooled statistics use 2008--2019 and 2021--2025. All income values are real hourly labor income in constant 2025 pesos and use GEIH expansion weights. Premiums are percent differences in weighted means relative to solo workers within each column.
  \end{minipage}
\end{table}


\begin{table}[htbp]
  \centering
  \caption{Worker composition by firm size: pooled sample and 2025}
  \label{tab:descriptive-worker-characteristics-comparison}
  \small
  \begin{tabular}{lrrrrrr}
    \toprule
    & \multicolumn{2}{c}{Women} & \multicolumn{2}{c}{Formal} & \multicolumn{2}{c}{Higher education} \\
    \cmidrule(lr){2-3} \cmidrule(lr){4-5} \cmidrule(lr){6-7}
    Firm size & Pooled & 2025 & Pooled & 2025 & Pooled & 2025 \\
    \midrule
    Solo & 43.6\% & 41.3\% & 6.7\% & 8.3\% & 13.1\% & 16.7\% \\
    2-3 & 33.2\% & 34.6\% & 11.3\% & 13.0\% & 14.2\% & 17.0\% \\
    4-5 & 32.5\% & 36.5\% & 22.1\% & 24.2\% & 18.3\% & 20.2\% \\
    6-10 & 36.4\% & 40.5\% & 41.6\% & 46.6\% & 26.1\% & 29.9\% \\
    11-19 & 39.2\% & 40.4\% & 60.7\% & 64.4\% & 35.0\% & 38.8\% \\
    20-30 & 40.6\% & 42.7\% & 74.6\% & 81.2\% & 40.4\% & 45.4\% \\
    31-50 & 40.3\% & 41.2\% & 84.2\% & 90.4\% & 44.2\% & 51.0\% \\
    51-100 & 40.8\% & 43.7\% & 89.5\% & 94.3\% & 47.1\% & 53.4\% \\
    101+ & 44.8\% & 45.8\% & 94.8\% & 98.0\% & 58.4\% & 62.9\% \\
    \bottomrule
  \end{tabular}
  \vspace{0.3em}
  \begin{minipage}{0.95\textwidth}
  \footnotesize
  Notes: Pooled statistics use 2008--2019 and 2021--2025. All statistics use GEIH expansion weights. Higher education corresponds to workers classified as superior or university educated.
  \end{minipage}
\end{table}


\begin{table}[htbp]
  \centering
  \caption{Real hourly labor income by firm size, pooled years}
  \label{tab:descriptive-income-size}
  \small
  \begin{tabular}{lrrrrrr}
    \toprule
    Firm size & Mean & P25 & Median & P75 & S.D. & Raw premium \\
    \midrule
    Solo & 5,690 & 2,366 & 3,943 & 6,169 & 11,205 & 0.0\% \\
    2-3 & 6,464 & 3,067 & 4,784 & 6,779 & 12,543 & 13.6\% \\
    4-5 & 7,444 & 3,943 & 5,430 & 7,163 & 21,476 & 30.8\% \\
    6-10 & 8,403 & 4,647 & 5,922 & 7,854 & 16,634 & 47.7\% \\
    11-19 & 9,746 & 5,204 & 6,383 & 8,774 & 17,881 & 71.3\% \\
    20-30 & 9,972 & 5,468 & 6,669 & 9,341 & 21,069 & 75.3\% \\
    31-50 & 10,744 & 5,714 & 6,939 & 10,321 & 16,435 & 88.8\% \\
    51-100 & 11,612 & 5,809 & 7,176 & 11,255 & 18,818 & 104.1\% \\
    101+ & 14,650 & 6,166 & 8,624 & 16,169 & 20,771 & 157.5\% \\
    \bottomrule
  \end{tabular}
  \vspace{0.3em}
  \begin{minipage}{0.95\textwidth}
  \footnotesize
  Notes: Incomes are pooled over 2008--2019 and 2021--2025 and expressed as real hourly labor income in constant 2025 pesos. All statistics are weighted by GEIH expansion weights. The raw premium is the percent difference in the weighted mean relative to solo workers.
  \end{minipage}
\end{table}


\begin{table}[htbp]
  \centering
  \caption{Worker characteristics by firm size, pooled years}
  \label{tab:descriptive-worker-characteristics}
  \small
  \begin{tabular}{lrrrr}
    \toprule
    Firm size & Women & Formal & Higher education & Mean log wage \\
    \midrule
    Solo & 49.7\% & 11.1\% & 19.9\% & 8.457 \\
    2-3 & 41.5\% & 17.1\% & 22.9\% & 8.596 \\
    4-5 & 40.2\% & 30.9\% & 26.7\% & 8.744 \\
    6-10 & 42.1\% & 50.0\% & 33.5\% & 8.847 \\
    11-19 & 43.8\% & 67.0\% & 41.6\% & 8.949 \\
    20-30 & 44.0\% & 78.8\% & 45.8\% & 8.986 \\
    31-50 & 43.2\% & 86.8\% & 48.9\% & 9.049 \\
    51-100 & 43.7\% & 91.2\% & 52.3\% & 9.114 \\
    101+ & 45.9\% & 95.6\% & 62.2\% & 9.300 \\
    \bottomrule
  \end{tabular}
  \vspace{0.3em}
  \begin{minipage}{0.95\textwidth}
  \footnotesize
  Notes: Statistics are pooled over 2008--2019 and 2021--2025 and weighted by GEIH expansion weights. Higher education corresponds to workers classified as superior or university educated.
  \end{minipage}
\end{table}


\begin{table}[htbp]
  \centering
  \caption{Real hourly labor income by firm size, 2025}
  \label{tab:descriptive-income-size-2025}
  \small
  \begin{tabular}{lrrrrrr}
    \toprule
    Firm size & Mean & P25 & Median & P75 & S.D. & Raw premium \\
    \midrule
    Solo & 8,351 & 3,606 & 5,769 & 8,942 & 11,807 & 0.0\% \\
    2-3 & 8,830 & 4,615 & 6,522 & 8,654 & 13,608 & 5.7\% \\
    4-5 & 9,503 & 5,440 & 6,923 & 8,654 & 16,211 & 13.8\% \\
    6-10 & 10,609 & 6,393 & 7,466 & 9,532 & 17,310 & 27.0\% \\
    11-19 & 12,346 & 6,844 & 7,867 & 11,276 & 18,438 & 47.8\% \\
    20-30 & 11,997 & 7,139 & 8,145 & 11,538 & 13,414 & 43.6\% \\
    31-50 & 13,428 & 7,141 & 8,515 & 13,516 & 17,573 & 60.8\% \\
    51-100 & 14,481 & 7,212 & 8,897 & 14,685 & 18,799 & 73.4\% \\
    101+ & 18,502 & 7,615 & 11,037 & 20,769 & 21,986 & 121.5\% \\
    \bottomrule
  \end{tabular}
  \vspace{0.3em}
  \begin{minipage}{0.95\textwidth}
  \footnotesize
  Notes: Incomes use only 2025 observations and are expressed as real hourly labor income in constant 2025 pesos. All statistics are weighted by GEIH expansion weights. The raw premium is the percent difference in the weighted mean relative to solo workers in 2025.
  \end{minipage}
\end{table}


\begin{table}[htbp]
  \centering
  \caption{Worker characteristics by firm size, 2025}
  \label{tab:descriptive-worker-characteristics-2025}
  \small
  \begin{tabular}{lrrrr}
    \toprule
    Firm size & Women & Formal & Higher education & Mean log wage \\
    \midrule
    Solo & 46.2\% & 14.1\% & 25.2\% & 8.639 \\
    2-3 & 42.4\% & 20.1\% & 26.1\% & 8.769 \\
    4-5 & 42.6\% & 32.2\% & 28.3\% & 8.879 \\
    6-10 & 44.0\% & 54.4\% & 36.3\% & 9.009 \\
    11-19 & 44.2\% & 70.4\% & 46.1\% & 9.128 \\
    20-30 & 45.1\% & 84.2\% & 49.5\% & 9.168 \\
    31-50 & 44.2\% & 91.9\% & 56.3\% & 9.254 \\
    51-100 & 46.8\% & 94.8\% & 58.4\% & 9.311 \\
    101+ & 47.5\% & 98.3\% & 67.6\% & 9.511 \\
    \bottomrule
  \end{tabular}
  \vspace{0.3em}
  \begin{minipage}{0.95\textwidth}
  \footnotesize
  Notes: Statistics use only 2025 observations and are weighted by GEIH expansion weights. Higher education corresponds to workers classified as superior or university educated.
  \end{minipage}
\end{table}


\begin{table}[!htbp]
  \centering
  \caption{Growth in mean real hourly labor income by firm size, 2008--2025}
  \label{tab:descriptive-wage-growth-2008-2025}
  \small
  \begin{tabular}{lrrr}
    \toprule
    Firm size & 2008 mean & 2025 mean & Change \\
    \midrule
    All firm sizes & 8,736 & 12,647 & 44.8\% \\
    Solo & 6,100 & 8,351 & 36.9\% \\
    2-3 & 7,223 & 8,830 & 22.2\% \\
    4-5 & 8,291 & 9,503 & 14.6\% \\
    6-10 & 8,249 & 10,609 & 28.6\% \\
    11-19 & 9,193 & 12,346 & 34.3\% \\
    20-30 & 9,347 & 11,997 & 28.3\% \\
    31-50 & 10,960 & 13,428 & 22.5\% \\
    51-100 & 11,193 & 14,481 & 29.4\% \\
    101+ & 13,369 & 18,502 & 38.4\% \\
    \bottomrule
  \end{tabular}
  \vspace{0.3em}
  \begin{minipage}{0.95\textwidth}
  \footnotesize
  Notes: Values are weighted mean real hourly labor income in constant 2025 pesos. Change is the percent change in the weighted mean between 2008 and 2025.
  \end{minipage}
\end{table}


\begin{table}[!htbp]
  \centering
  \caption{Mean real hourly labor income by firm size and worker group, 2008 and 2025}
  \label{tab:descriptive-wage-groups-mean-2008-2025}
  \small
  \begin{tabular}{lrrrrr}
    \toprule
    Firm size & All & Men & Women & Formal & Informal \\
    \midrule
    \multicolumn{6}{l}{\textit{Panel A: 2008}} \\
    All firm sizes & 8,736 & 9,087 & 8,264 & 12,476 & 6,367 \\
    Solo & 6,100 & 6,790 & 5,357 & 12,658 & 5,626 \\
    2-3 & 7,223 & 7,355 & 6,980 & 11,377 & 6,674 \\
    4-5 & 8,291 & 8,680 & 7,585 & 11,831 & 7,204 \\
    6-10 & 8,249 & 8,788 & 7,399 & 10,061 & 6,890 \\
    11-19 & 9,193 & 9,807 & 8,372 & 10,187 & 7,796 \\
    20-30 & 9,347 & 9,406 & 9,259 & 10,256 & 7,261 \\
    31-50 & 10,960 & 10,654 & 11,401 & 10,719 & 11,901 \\
    51-100 & 11,193 & 10,589 & 12,111 & 11,911 & 7,075 \\
    101+ & 13,369 & 13,447 & 13,270 & 13,728 & 9,622 \\
    \addlinespace
    \multicolumn{6}{l}{\textit{Panel B: 2025}} \\
    All firm sizes & 12,647 & 12,849 & 12,408 & 16,571 & 7,401 \\
    Solo & 8,351 & 8,749 & 7,889 & 17,217 & 6,894 \\
    2-3 & 8,830 & 8,828 & 8,832 & 14,077 & 7,511 \\
    4-5 & 9,503 & 9,561 & 9,424 & 13,692 & 7,517 \\
    6-10 & 10,609 & 10,824 & 10,335 & 12,844 & 7,947 \\
    11-19 & 12,346 & 12,833 & 11,733 & 13,656 & 9,227 \\
    20-30 & 11,997 & 11,957 & 12,045 & 12,521 & 9,206 \\
    31-50 & 13,428 & 13,484 & 13,358 & 13,879 & 8,337 \\
    51-100 & 14,481 & 14,992 & 13,900 & 14,598 & 12,327 \\
    101+ & 18,502 & 18,869 & 18,096 & 18,548 & 15,843 \\
    \bottomrule
  \end{tabular}
  \vspace{0.3em}
  \begin{minipage}{0.95\textwidth}
  \footnotesize
  Notes: Values are weighted mean real hourly labor income in constant 2025 pesos. The first row pools all firm-size categories. Men and women are defined from the worker's reported sex; formal and informal are defined using the baseline formality indicator.
  \end{minipage}
\end{table}


\begin{table}[!htbp]
  \centering
  \caption{Median real hourly labor income by firm size and worker group, 2008 and 2025}
  \label{tab:descriptive-wage-groups-median-2008-2025}
  \small
  \begin{tabular}{lrrrrr}
    \toprule
    Firm size & All & Men & Women & Formal & Informal \\
    \midrule
    \multicolumn{6}{l}{\textit{Panel A: 2008}} \\
    All firm sizes & 5,128 & 5,244 & 5,034 & 7,132 & 4,195 \\
    Solo & 4,027 & 4,195 & 3,596 & 6,616 & 3,776 \\
    2-3 & 4,489 & 4,720 & 4,195 & 5,873 & 4,195 \\
    4-5 & 4,840 & 4,894 & 4,840 & 5,422 & 4,661 \\
    6-10 & 5,244 & 5,244 & 5,244 & 5,802 & 4,840 \\
    11-19 & 5,559 & 5,454 & 5,664 & 6,162 & 5,003 \\
    20-30 & 5,929 & 5,768 & 6,293 & 6,413 & 5,034 \\
    31-50 & 6,075 & 5,873 & 6,293 & 6,293 & 4,952 \\
    51-100 & 6,293 & 6,293 & 6,678 & 6,670 & 5,034 \\
    101+ & 8,055 & 7,866 & 8,390 & 8,390 & 5,387 \\
    \addlinespace
    \multicolumn{6}{l}{\textit{Panel B: 2025}} \\
    All firm sizes & 7,805 & 7,805 & 7,821 & 9,615 & 5,769 \\
    Solo & 5,769 & 6,000 & 5,538 & 10,256 & 5,385 \\
    2-3 & 6,522 & 6,593 & 6,410 & 8,395 & 5,769 \\
    4-5 & 6,923 & 6,923 & 6,923 & 8,145 & 6,250 \\
    6-10 & 7,466 & 7,466 & 7,466 & 8,077 & 6,731 \\
    11-19 & 7,867 & 7,867 & 7,867 & 8,308 & 6,923 \\
    20-30 & 8,145 & 8,027 & 8,212 & 8,308 & 7,023 \\
    31-50 & 8,515 & 8,392 & 8,654 & 8,654 & 7,133 \\
    51-100 & 8,897 & 8,654 & 9,135 & 9,030 & 7,692 \\
    101+ & 11,037 & 10,714 & 11,538 & 11,037 & 7,692 \\
    \bottomrule
  \end{tabular}
  \vspace{0.3em}
  \begin{minipage}{0.95\textwidth}
  \footnotesize
  Notes: Values are weighted median real hourly labor income in constant 2025 pesos. The first row pools all firm-size categories. Men and women are defined from the worker's reported sex; formal and informal are defined using the baseline formality indicator.
  \end{minipage}
\end{table}



\subsection{Additional Descriptive Figures}

Figures~\ref{fig:descriptive-employment-share-time},
\ref{fig:descriptive-raw-premium-time}, and
\ref{fig:descriptive-worker-composition} report the full time path of
employment shares, raw firm-size wage premia, and worker composition.
The time-series figures show persistence rather than a one-year anomaly:
employment remains concentrated at the bottom, raw wage premia remain positive,
and formality and higher education remain strongly increasing in firm size.

\begin{figure}[htbp]
  \centering
  \includegraphics[width=0.9\textwidth]{fig61_\figurelang.png}
  \caption{Employment share by firm size over time}
  \label{fig:descriptive-employment-share-time}
\end{figure}

\begin{figure}[htbp]
  \centering
  \includegraphics[width=0.9\textwidth]{fig62_\figurelang.png}
  \caption{Raw firm-size wage premium over time, relative to solo workers}
  \label{fig:descriptive-raw-premium-time}
\end{figure}

\begin{figure}[htbp]
  \centering
  \includegraphics[width=0.9\textwidth]{fig63_\figurelang.png}
  \caption{Worker composition by firm size}
  \label{fig:descriptive-worker-composition}
\end{figure}

\clearpage

\subsection{Additional Regression Tables and Figures}
\label{app:additional-regression-figures}

Figure~\ref{fig:firm-size-beta-coefficients} reports the log-point coefficients
from the full specification, and Figure~\ref{fig:firm-size-premium-trend-test}
visualizes the corresponding trend-test estimates.
Together, these figures show that the conditional firm-size coefficients are
positive in the full specification, while the estimated time trends are positive
but too imprecise to support a strong claim that the premium has systematically
steepened.

Table~\ref{tab:firm-size-local-geography-robustness} reports a local-market
robustness exercise. Because the harmonized file retains department/Bogot\'a
rather than a consistent city identifier, the table uses department as the local
geography. The most demanding column absorbs sector-department-year fixed
effects, so the firm-size coefficients are identified within the same sector,
department, and year. The large-firm coefficient remains positive in this
specification: workers in firms with 101 or more employees earn about 31 percent
more than comparable solo workers within the same sector-department-year cell.

\begin{table}[htbp]
  \centering
  \caption{Firm-size wage premium robustness with local sector-year fixed effects}
  \label{tab:firm-size-local-geography-robustness}
  \small
  \begin{tabular}{lccc}
    \toprule
    & \multicolumn{3}{c}{Dependent variable: log real hourly labor income} \\
    \cmidrule(lr){2-4}
    & (1) & (2) & (3) \\
    \midrule
    Firm size: 2-3 & 0.145*** & 0.126*** & 0.127*** \\
     & (0.014) & (0.014) & (0.015) \\
    Firm size: 4-5 & 0.234*** & 0.216*** & 0.215*** \\
     & (0.020) & (0.020) & (0.021) \\
    Firm size: 6-10 & 0.234*** & 0.215*** & 0.215*** \\
     & (0.023) & (0.023) & (0.024) \\
    Firm size: 11-19 & 0.224*** & 0.209*** & 0.209*** \\
     & (0.022) & (0.022) & (0.023) \\
    Firm size: 20-30 & 0.190*** & 0.177*** & 0.177*** \\
     & (0.022) & (0.022) & (0.023) \\
    Firm size: 31-50 & 0.195*** & 0.184*** & 0.183*** \\
     & (0.022) & (0.023) & (0.023) \\
    Firm size: 51-100 & 0.215*** & 0.203*** & 0.203*** \\
     & (0.022) & (0.022) & (0.023) \\
    Firm size: 101+ & 0.282*** & 0.272*** & 0.268*** \\
     & (0.030) & (0.032) & (0.032) \\
    \midrule
    Worker and formality controls & Yes & Yes & Yes \\
    Fixed effects & $s\times t$ & $s\times t$, $d\times t$ & $s\times d\times t$ \\
    Observations & 3,852,322 & 3,852,322 & 3,852,211 \\
    $R^2$ & 0.362 & 0.374 & 0.380 \\
    \bottomrule
  \end{tabular}
  \vspace{0.3em}
  \begin{minipage}{0.95\textwidth}
  \footnotesize
  Notes: The omitted category is solo workers. The sample is restricted to observations with nonmissing department codes. Worker controls include a female-worker indicator, age, age squared, education dummies, and labor formality. The local geography $d$ is department/Bogot\'a, because the harmonized file does not retain a consistent city or municipality identifier. Column (3) is the most stringent specification and absorbs sector-department-year fixed effects, so identification comes from wage differences across firm-size categories within the same sector, department, and year. Standard errors are clustered by sector-department cells. Significance levels: * $p<0.10$, ** $p<0.05$, *** $p<0.01$.
  \end{minipage}
\end{table}


\begin{figure}[htbp]
  \centering
  \includegraphics[width=0.9\textwidth]{fig58_\figurelang.png}
  \caption{Firm-size coefficients in the full specification}
  \label{fig:firm-size-beta-coefficients}
\end{figure}

\begin{figure}[htbp]
  \centering
  \includegraphics[width=0.9\textwidth]{fig60_\figurelang.png}
  \caption{Trend test for the firm-size wage premium}
  \label{fig:firm-size-premium-trend-test}
\end{figure}
